\chapter{릴레이식의 박자: 순서기}\label{ch:relay-sequencer}

기어식에서는 주기 축이 돌며 슬롯 원판이 시각을 냈다(\ref{sec:t-gear}절). 릴레이식에는 도는 축이 없으므로 시각을
따로 만들어야 한다. 이 장은 \ref{ch:relay-timing}장의 시간 규칙을 릴레이식이 어떻게 지키는지 본다.

\section{릴레이식: 555와 링 계수기가 박자를 친다}\label{sec:t-sequencer}

릴레이식 뱅크의 \term{순서기}{sequencer} 보드는 위치마다 \code{STROBE},
\code{JUDGE}, \code{CLEAR}, 그리고 계수기의 \code{CK}/\code{NCK}를 L1--L9에 맞춰 낸다. 정지와 재출발,
\code{HOME}(계수기 프리셋), 마지막 위치(\code{END}), 대기도 이 보드가 처리한다. 릴레이 50개가 있다(폴더
\code{hardware/relay/sequencer}).

\Why 순서·조건·잠금은 모두 릴레이 접점으로 만들었다. 그런데 \emph{박자}, 곧 ``6.3~ms마다 다음 단계''라는 시간
기준 하나만은 전자식이다. 이유는 \ref{sec:t-bounds}절의 하한에 있다. 규칙이 하한이면 시간 기준의 \emph{최소}를
보증할 수 있어야 한다. 릴레이로 지연을 만들려면 코일에 커패시터를 달아 천천히 충전되게 하는데, 릴레이가 붙는
전압은 ``정격의 80\% 이하''처럼 위쪽 한계만 알려져 있다. 실제로 50\%에서 붙는 개체와 80\%에서 붙는 개체는 같은
RC에서 지연이 두 배 넘게 다르다. 하한을 보증할 수 없다. 릴레이를 사슬로 이어 박자를 만들면 단마다 흩어짐이
쌓여 한 주기가 70--80~ms로 늘어난다. 그래서 전자 부품 두 개로 박자를 만들고, 그 박자가 정하는 것은 ``언제''뿐이며
``진행할지''는 여전히 릴레이가 정한다. 기어식에서 슬롯 원판이 하는 일과 같다(\ref{sec:t-gear}절).

\General 시간 기준은 입문 전자 회로의 단골 둘이다.
\begin{itemize}
\item \textbf{555 타이머}는 커패시터 하나를 저항을 거쳐 충전하다가 전원의 2/3에 이르면 방전으로 바꾸고, 1/3까지
내려가면 다시 충전으로 바꾸는 칩이다. 출력은 충전 중에 높고 방전 중에 낮은 네모파가 된다. 한 주기는
\[
  T = 0.693\,(R_1 + 2R_2)\,C
\]
이다. 충전은 $R_1 + R_2$를 지나고 방전은 $R_2$만 지나기 때문에 $R_2$가 두 번 들어간다.
\item \textbf{74HC4017}은 클럭이 올 때마다 출력 Q0, Q1, \ldots, Q9 가운데 \emph{하나만} 차례로 켜는 칩이다(10단
링 계수기). 흐르는 불빛 장식을 만드는 입문 예제에 꼭 나온다.
\end{itemize}

\Ours 순서기는 LMC555($R_1$ = 10~k$\Omega$, $R_2$ = 150~k$\Omega$ + 트리머 100~k$\Omega$, $C$ = 22~nF)가 6.3~ms마다
4017을 한 칸 민다. 열 칸이 한 위치이므로 한 주기는 63~ms다. 4017의 출력 몇 개를 다이오드 OR로 묶어
\term{창}{window}을 만들고, 창이 고측 구동기(\ref{sec:el-mosfet}절)를 거쳐 출력 릴레이의 코일을 먹인다(그림~\ref{fig:t-ring}).

\begin{figure}[ht]
\centering
\begin{tikzpicture}[x=1.15cm,y=0.62cm]
  \foreach \q in {0,...,9} {
    \draw[cgray!60] (\q,0.6) -- (\q,-5.6);
    \node[font=\scriptsize\ttfamily] at (\q+0.5,1.0) {Q\q};
  }
  \draw[cgray!60] (10,0.6) -- (10,-5.6);
  \newcommand{\seqwin}[5]{\fill[#4] (#1,#3-0.35) rectangle (#2,#3+0.35); \node[font=\scriptsize,anchor=east] at (-0.15,#3) {#5};}
  \seqwin{0}{2}{0}{cbluebg}{\code{STROBE} (Q0--Q1)}
  \seqwin{3}{5}{-1.2}{cgreenbg}{\code{JUDGE} (Q3--Q4)}
  \seqwin{6}{9}{-2.4}{corangebg}{\code{CLEAR} (Q6--Q8)}
  \seqwin{6}{7}{-3.6}{cyellowbg}{\code{CK} 결정 (Q6)}
  \seqwin{7}{9}{-4.8}{credbg}{\code{CK} (Q7--Q8)}
  \node[font=\scriptsize,anchor=west] at (9.05,-4.8) {Q9: 주차};
  \draw[<->] (0,-6.2) -- node[below,font=\scriptsize] {한 틱 $T$ = 6.3 ms (관찰 모드 27 ms)} (1,-6.2);
  \draw[<->] (0,-7.4) -- node[below,font=\scriptsize] {한 위치 = 10틱 = 63 ms (관찰 모드 270 ms)} (10,-7.4);
\end{tikzpicture}
\caption{순서기의 링 상태와 창. 창이 바로 출력이 되지는 않는다. 창과 출력 릴레이의 코일 사이에 상태 릴레이의
접점(순환 허가, \code{CK} 결정 등)이 직렬로 들어가, 접점이 닫혀 있어야만 창이 코일에 닿는다.}
\label{fig:t-ring}
\end{figure}

왜 63~ms인가. 출력이 릴레이 접점이므로 순서기 \emph{자신의} 출력 릴레이도 붙고 놓이는 데 시간이 든다.
\code{STROBE}의 핀 폭은 구동 창의 폭에서 (붙는 시간 $-$ 놓이는 시간)만큼 깎여 최악 3.5--4.1~ms 줄어든다. 그래서
L1의 7.5~ms를 지키려면 구동 창이 11.6~ms 이상, 곧 두 틱이 5.8~ms 이상이어야 하고, 여유를 두어 6.3~ms로 정했다.
시뮬레이션이 릴레이 시간을 가장 나쁘게 넣어 찾은 \code{STROBE}의 최소 폭은 8.23~ms(여유 0.73~ms)였다.

\code{CK}를 낼지는 Q6에서 정한다. 그때 \code{RUN}(정지가 아님)이 높으면 \code{CK} 결정 릴레이가 붙어 Q6--Q8 동안
스스로를 붙들고, Q7--Q8에 \code{CK}가 나간다. 정지가 나서 \code{RUN}이 내려가면, 릴레이 시간이 가장 나쁜
조합에서도 Q6이 시작된 뒤 0.6~ms 안에 \code{RUN}을 따르는 릴레이가 열려, 붙는 데 2.5~ms 이상 걸리는 결정
릴레이의 동작이 취소된다(시뮬레이션으로 확인). 이 순서가 중요하다. 계수기는 내려 셀 수 없으므로 정지한 위치를 한 번 지나치면 되돌아올
방법이 없다.

\subsection{관찰 모드: 왜 커패시터가 아니라 저항을 바꾸는가}

전시에서는 관람객이 릴레이가 넘어가는 것을 눈으로 따라갈 수 있게 느린 \term{관찰 모드}{observation mode}가
필요하다. 순서기는 555의 $R_2$에 680~k$\Omega$을 직렬로 넣어 틱을 27.0~ms로 늘린다(초당 3.7 위치). 평소에는
릴레이 접점 두 개가 직렬로 이 저항을 단락해 두고, 관찰 쪽으로 스위치를 돌리면 단락이 열린다.

\Why 속도를 바꾸는 순간은 운전 중 아무 때나다. 관람객이 손잡이를 돌리기 때문이다. 하한 규칙에서 위험한 것은
틱이 \emph{짧아지는} 것 하나뿐이므로, 바꾸는 순간이 걸린 틱도 빠른 틱보다 짧아지면 안 된다(그림~\ref{fig:t-555}).

\begin{itemize}
\item \textbf{저항을 바꾸면} 커패시터의 전압은 그대로 이어지고 충전·방전의 빠르기만 느려진다. 전압이 빠른
모드보다 빨리 움직일 수 없으므로 어느 틱도 빠른 틱보다 짧아지지 않는다. 링은 한 번 Q9를 떠나면 반드시 Q9까지
돌므로 주기가 중간에 끊기지도 않는다.
\item \textbf{커패시터를 더하면} 충전된 커패시터와 빈 커패시터가 전하를 나눠 전압이 순간 떨어진다. 방전 중이었다면
곧바로 1/3 아래로 내려가 555가 충전으로 넘어간다. 그 틱이 반쯤에서 잘려 \code{STROBE}나 \code{JUDGE}의 폭이
규칙보다 짧아질 수 있다.
\item \textbf{분주기}(4017을 하나 더 두어 느린 클럭을 고르는 것)도 두 클럭을 비동기로 바꾸는 순간 틱이 임의로
짧아진다. 막으려면 동기화 회로가 필요해 전자 부품이 는다.
\end{itemize}

\begin{figure}[ht]
\centering
\begin{tikzpicture}[x=1cm,y=2.2cm]
  \foreach \dx/\ttl in {0/저항을 바꿈,7.6/커패시터를 더함} {
    \begin{scope}[xshift=\dx cm]
      \node[font=\small\bfseries] at (3,1.25) {\ttl};
      \draw[->] (0,0) -- (6.4,0) node[right,font=\scriptsize] {$t$};
      \draw[->] (0,0) -- (0,1.1) node[above,font=\scriptsize] {$v_C$};
      \draw[cgray,dashed] (0,1/3) -- (6.2,1/3) node[right,font=\tiny] {1/3};
      \draw[cgray,dashed] (0,2/3) -- (6.2,2/3) node[right,font=\tiny] {2/3};
      % two normal ticks: charge (tau 1) 1/3->2/3 takes ln2 = 0.693, discharge same
      \draw[cblue,line width=0.9pt,domain=0:0.693,samples=20] plot (\x,{1-(2/3)*exp(-\x)});
      \draw[cblue,line width=0.9pt,domain=0:0.693,samples=20] plot (\x+0.693,{(2/3)*exp(-\x)});
      \draw[cblue,line width=0.9pt,domain=0:0.693,samples=20] plot (\x+1.386,{1-(2/3)*exp(-\x)});
      \draw[cred] (2.379,-0.05) -- (2.379,1.05);
      \node[font=\tiny,text=cred,anchor=south] at (2.379,1.05) {바꾸는 순간};
    \end{scope}}
  % left: discharge starts at 2.079 with tau 1 for 0.3, then tau 4
  \draw[cblue,line width=0.9pt,domain=0:0.3,samples=10] plot (\x+2.079,{(2/3)*exp(-\x)});
  \draw[corange,line width=0.9pt,domain=0:1.57,samples=30] plot (\x+2.379,{(2/3)*exp(-0.3)*exp(-\x/4)});
  \draw[corange,line width=0.9pt,domain=0:0.9,samples=20] plot (\x+3.949,{1-(2/3)*exp(-\x/4)});
  \node[note,anchor=west,align=left,text=corange] at (2.5,0.15) {전압은 이어지고\\내려가는 빠르기만 느려진다};
  % right: discharge 0.3 then voltage halves (equal empty capacitor added) -> below 1/3 -> charging restarts
  \begin{scope}[xshift=7.6cm]
    \draw[cblue,line width=0.9pt,domain=0:0.3,samples=10] plot (\x+2.079,{(2/3)*exp(-\x)});
    \draw[corange,line width=0.9pt] (2.379,{(2/3)*exp(-0.3)}) -- (2.379,{(1/3)*exp(-0.3)});
    \draw[corange,line width=0.9pt,domain=0:1.6,samples=30] plot (\x+2.379,{1-(1-(1/3)*exp(-0.3))*exp(-\x/2)});
    \node[note,anchor=west,align=left,text=cred] at (2.6,0.12) {1/3 아래로 떨어져\\틱이 반쯤에서 잘린다};
  \end{scope}
\end{tikzpicture}
\caption{555 커패시터 전압의 모양(개념도). 왼쪽은 방전 도중에 $R_2$를 늘린 경우로, 그 틱은 오히려 길어진다.
오른쪽은 같은 순간에 빈 커패시터를 병렬로 더한 경우로, 전압이 순간 반으로 떨어져 555가 곧바로 충전으로
넘어간다. 그 틱은 짧아진다.}
\label{fig:t-555}
\end{figure}

\Ours 순서기의 시뮬레이터는 555를 넷리스트의 RC로 흉내 낸다. 그 순간 닫혀 있는 접점으로 이어진 저항의 합을
$R_2$로 쓰고, 저항이 바뀌면 커패시터 전압을 이어서 계산한다(\ref{sec:v-time}절). 운전 중 선택 스위치를 무작위
시각에 36번씩 뒤집어도 가장 짧은 555 주기가 빠른 틱의 0.997배였다(넣어 둔 흔들림 $\pm$0.3\% 안). 반대로
``관찰 저항이 들어가지 않는'' 고장을 일부러 넣은 회로는 주기가 59.9~ms로 나와 검증기에 걸린다.

\subsection{한 걸음 버튼}

관찰 모드보다 더 느린 것이 \term{한 걸음}{single step} 버튼이다. 한 번 누르면 정확히 한 위치만 진행한다.

\Why ``정확히 한 바퀴''는 이미 4017이 보장한다. 4017은 Q9에서만 서도록(주차) 배선되어 있어, 링이 한 번 Q9를
떠나면 \code{STROBE}, \code{JUDGE}, \code{CLEAR}, \code{CK}를 모두 낸 뒤 Q9로 돌아와 선다. 그래서 버튼 회로가 할
일은 ``한 바퀴를 시작하게 하고, 시작했으면 다음 바퀴는 막는 것''뿐이다. 이것을 릴레이 두 개의 래치로 만든다.
\begin{itemize}
\item \textbf{받아들임} 릴레이(\code{KSA}): 링이 주차해 있을 때 버튼이 눌리면 선다. 링을 출발시키는 조건(GO)에
들어간다.
\item \textbf{주기} 릴레이(\code{KSC}): 받아들인 한 걸음의 링이 Q0에 들어서면 선다. 서자마자 그 NC 접점이 한
걸음의 GO를 끊는다. 이번 바퀴는 이미 출발했으므로 계속 돌고, 다음 바퀴는 출발하지 않는다.
\end{itemize}
버튼을 계속 누르고 있으면 두 릴레이가 서로를 붙들어 다음 바퀴가 없다(반복하지 않는다). 버튼을 일찍 놓아도
받아들임 릴레이는 링이 출발한 동안 다른 접점으로 붙들려 한 바퀴를 마친다.

\General 컴퓨터의 키보드나 마우스 버튼도 ``누를 때마다 정확히 한 번''이 필요하다. 접점은 닫힐 때 수 ms 동안
여러 번 튀기 때문에 그대로 세면 한 번 누름이 여러 번이 된다. 보통 펌웨어가 일정 시간 기다려 튐을 거른다.

\Ours 이 보드는 튐을 릴레이의 느림으로 거른다. 버튼을 따라가는 릴레이(\code{KSP})는 붙는 데 2.5~ms 이상이
걸리므로 그보다 짧게 닫히는 튐에는 붙지 않고, 코일에 다이오드만 달아 놓이는 데 3~ms 이상이 걸리므로 그보다
짧게 열리는 튐에는 놓이지 않는다. 받아들임은 주차 중에만 되므로 바퀴가 도는 동안의 튐은 아무 일도 하지 않는다.
한 걸음 주기는 선택 스위치와 무관하게 늘 관찰 속도로 돈다. 한 걸음 중에는 \code{RUN}이 낮으므로 정지가 나면
\code{STOP\_LATCH}를 거치는 조금 느린 길로 \code{CK}를 막아야 하는데, 빠른 틱에서는 그 여유가 0.1~ms뿐이기
때문이다.


\section{정리}

\begin{itemize}
\item 릴레이식 순서기는 555와 4017로 박자만 만들고 나머지는 릴레이로 정한다. 관찰 모드는 커패시터 전압이 이어지는
저항 쪽을 바꿔 어떤 틱도 짧아지지 않게 한다. 한 걸음 버튼은 ``시작은 허락하고 다음 바퀴는 막는'' 래치다.
\item 시간 기준이 아무리 흔들려도 틱이 짧아지지 않으므로 L1--L9가 성립하고, 순서와 조건은 릴레이 접점이 지킨다.
\end{itemize}
